\documentclass[11pt]{article}
\usepackage[margin=0.8in]{geometry}
\usepackage{amsmath,amssymb,amsthm}
\usepackage{enumitem}
\usepackage{longtable}
\usepackage{array}
\usepackage{xcolor}
\usepackage{hyperref}
\usepackage{fancyhdr}
\usepackage{titlesec}
\usepackage{tcolorbox}
\usepackage{multicol}
\usepackage{booktabs}
\hypersetup{colorlinks=true, linkcolor=blue!50!black, urlcolor=blue!50!black}
\pagestyle{fancy}
\fancyhf{}
\rhead{\thepage}
\lhead{Engineering Math to Pure Math}
\setlength{\parskip}{0.55em}
\setlength{\parindent}{0pt}
\setlist[itemize]{topsep=0.2em,itemsep=0.2em}
\setlist[enumerate]{topsep=0.2em,itemsep=0.25em}
\newtcolorbox{focusbox}{colback=blue!3!white,colframe=blue!50!black,title=Focus}
\newtcolorbox{intuitionbox}{colback=green!3!white,colframe=green!40!black,title=Intuition}
\newtcolorbox{formalbox}{colback=gray!4!white,colframe=gray!70!black,title=Formal}
\newtcolorbox{walkbox}{colback=orange!5!white,colframe=orange!60!black,title=Walking question}
\newcommand{\N}{\mathbb{N}}
\newcommand{\Z}{\mathbb{Z}}
\newcommand{\Q}{\mathbb{Q}}
\newcommand{\R}{\mathbb{R}}
\newcommand{\C}{\mathbb{C}}
\newcommand{\eps}{\varepsilon}
\newcommand{\pow}{\mathcal{P}}
\theoremstyle{definition}
\newtheorem{definition}{Definition}
\newtheorem{example}{Example}
\newtheorem{exercise}{Exercise}
\theoremstyle{plain}
\newtheorem{theorem}{Theorem}
\newtheorem{lemma}{Lemma}
\theoremstyle{remark}
\newtheorem{remark}{Remark}

\title{\textbf{Proofs From Zero to Velleman}\\A Conversational Reader for Building Mathematical Maturity}
\author{Engineering Math to Pure Math}
\date{May 2026}
\begin{document}
\maketitle
\tableofcontents
\newpage

\section*{Preface: What This Reader Is}
\addcontentsline{toc}{section}{Preface: What This Reader Is}

This is the full self-contained reader rather than a syllabus or summary checklist: it works through the definitions, examples, and proof techniques directly.

The goal is simple to say and serious to do: take someone who knows useful engineering mathematics and help them become comfortable with proof-based pure mathematics. That means learning to read slowly, learning to use definitions as tools, learning to make examples and counterexamples, learning to write proofs that survive inspection, and learning why these little habits eventually open the doors to real analysis, abstract algebra, topology, differential geometry, and algebraic geometry.

The first major book is Velleman's \emph{How to Prove It}. This reader is meant to make the opening of that book feel familiar instead of foggy. We will start before Velleman starts: with mathematical sentences themselves.

The tone here is friendly, but the destination is rigor. A proof is not a vibe. It is not a pile of plausible sentences. It is a sequence of justified moves from assumptions to conclusion. But rigor does not have to arrive wearing a cold face. Mathematics was built by human beings trying to understand number, shape, motion, symmetry, infinity, and structure. The formal language came later because the ideas were powerful enough to need protection.

\begin{focusbox}
Read with a pencil. Do the small exercises. Say definitions aloud. If a sentence feels slippery, do not push through it. Stop and ask: what are the objects, what is being assumed, and what must be shown?
\end{focusbox}

\newpage
\section{The Change in Game}

Engineering mathematics often trains excellent reflexes. You see a differential equation; you try separation, integrating factors, Laplace transforms, eigenvalues. You see an integral; you look for substitution, parts, symmetry, a known transform. You see a matrix; you think rank, determinant, eigenvectors, conditioning.

That is real skill. Do not throw it away.

But pure mathematics asks a different question first. Instead of ``What method solves this?'', it asks:

\begin{quote}
What exactly is being claimed, and why must it follow from the assumptions?
\end{quote}

That one question changes the whole room.

\subsection{From calculation to justification}

Calculation is still present. Pure mathematicians compute all the time. But computation is no longer the whole story. A computation can suggest a pattern; a proof explains why the pattern cannot fail.

For example:
\[
1+3=4,\qquad 1+3+5=9,\qquad 1+3+5+7=16.
\]
You might guess:
\[
1+3+5+\cdots+(2n-1)=n^2.
\]

A calculator can check the first thousand cases. A proof covers all cases at once.

\subsection{The big shift}

The shift is from answer-getting to structure-seeing. You are learning to ask:

\begin{itemize}
\item What kind of object is this?
\item What is the exact definition?
\item What examples satisfy it?
\item What examples almost satisfy it but fail?
\item What is the theorem really promising?
\item What would break it?
\item Which assumptions are doing work?
\end{itemize}

\begin{walkbox}
On a walk, take one theorem you know from calculus, like ``continuous functions on closed intervals attain maximum values.'' Ask: which words are essential? Continuous? Closed? Bounded? Interval? What breaks if one word disappears?
\end{walkbox}

\newpage
\section{Statements, Predicates, and the First Filter}

Mathematics begins with sentences that can be judged true or false.

\begin{definition}
A \emph{statement} is a sentence that is either true or false.
\end{definition}

\begin{example}
``$5$ is prime'' is a statement. It is true.
\end{example}

\begin{example}
``$9$ is prime'' is a statement. It is false.
\end{example}

\begin{example}
``Find the derivative of $x^2$'' is not a statement. It is a command.
\end{example}

\subsection{Predicates}

\begin{definition}
A \emph{predicate} is a sentence with one or more variables that becomes a statement when the variables are assigned values from a domain.
\end{definition}

The sentence
\[
x^2>10
\]
is not true or false until $x$ is known. If $x=4$, true. If $x=2$, false. So $x^2>10$ is a predicate.

In software terms, a predicate is like a Boolean-valued function:

\[
P(x)=\text{``}x^2>10\text{''}.
\]

But mathematics demands something software often hides: the domain. Are we allowing $x\in\N$, $x\in\Z$, $x\in\R$, or $x\in\C$? The sentence may change character when the domain changes.

\subsection{Why this matters}

A beginner often reads a theorem too quickly. The theorem is a sentence with machinery inside it. Variables have domains. Quantifiers control choice. Hypotheses give permissions. Conclusions give obligations.

Before proving anything, learn to parse.

\subsection{Exercises}

For each item, decide whether it is a statement, predicate, command, or expression.

\begin{enumerate}
\item $2+2=4$.
\item $x+2=4$.
\item Solve $x+2=4$.
\item There are infinitely many prime numbers.
\item $n$ is divisible by $3$.
\item $\int_0^1 x^2\,dx$.
\item Every polynomial of odd degree with real coefficients has a real root.
\item Is every continuous function differentiable?
\end{enumerate}

\begin{focusbox}
Training target: never begin a proof until you know what the sentence is saying.
\end{focusbox}

\newpage
\section{Variables Are Not All the Same}

One source of mathematical fog is that every variable looks like a letter, but not every variable has the same role.

Some variables are fixed. Some are arbitrary. Some are chosen as witnesses. Some are temporary. Some are bound by quantifiers. Some are free.

\subsection{Fixed but unknown}

When a proof says:

\begin{quote}
Let $n\in\Z$.
\end{quote}

It means $n$ is now fixed for the discussion. We may not know which integer it is, but we are treating it as one particular integer.

\subsection{Arbitrary}

When a proof says:

\begin{quote}
Let $n\in\Z$ be arbitrary.
\end{quote}

It means we are not using any special property of $n$ except that it is an integer. If the argument works, it works for every integer.

This is how universal proofs begin.

\subsection{Witness}

When proving an existential statement, a variable may be chosen to prove existence.

Claim:
\[
\exists n\in\Z \text{ such that } n^2=81.
\]

Proof:

\begin{quote}
Take $n=9$. Then $n\in\Z$ and $n^2=81$.
\end{quote}

Here $9$ is the witness.

\subsection{Bound and free variables}

In
\[
\forall x\in\R,\ x^2\geq 0,
\]
the variable $x$ is bound by the quantifier $\forall$. It is local to that sentence.

In
\[
x^2\geq 0,
\]
the variable $x$ is free unless context has already assigned it a meaning.

This is close to scope in programming. A loop variable belongs to the loop. A local variable belongs to the function. A mathematical variable belongs to the sentence or proof context that introduced it.

\subsection{Exercises}

In each sentence, identify the role of the variable.

\begin{enumerate}
\item Let $x\in\R$ be arbitrary.
\item Choose $n=12$.
\item For every $a\in A$, there exists $b\in B$ such that $f(a)=b$.
\item Suppose $p$ is prime.
\item Let $\eps>0$.
\end{enumerate}

\newpage
\section{Quantifiers: The Tiny Symbols That Run the Machine}

The two most important symbols at the beginning are:
\[
\forall \quad \text{and} \quad \exists.
\]

\[
\forall = \text{for every}, \qquad \exists = \text{there exists}.
\]

Small symbols. Large consequences.

\subsection{Universal statements}

A universal statement says every object in a domain has some property:
\[
\forall x\in D,\ P(x).
\]

To prove it, take an arbitrary $x\in D$ and prove $P(x)$.

To disprove it, find one counterexample.

\subsection{Existential statements}

An existential statement says at least one object in a domain has some property:
\[
\exists x\in D \text{ such that } P(x).
\]

To prove it, provide a witness or prove one must exist.

To disprove it, prove every object in the domain fails the property.

\subsection{The order matters}

Compare:
\[
\forall x\in\R,\ \exists y\in\R \text{ such that } x+y=0.
\]
\[
\exists y\in\R \text{ such that } \forall x\in\R,\ x+y=0.
\]

The first is true. Each $x$ gets its own $y=-x$.

The second is false. It asks for one $y$ that cancels every $x$.

This is the sort of thing that makes proofs feel mysterious if you read too fast. The symbols are not decoration. They describe a game of choice.

\subsection{A game interpretation}

Think of $\forall$ as an adversary choosing an object. Think of $\exists$ as you responding with an object.

For
\[
\forall x\in\R,\ \exists y\in\R,\ x+y=0,
\]
the adversary hands you $x$. You respond with $y=-x$. You win every time.

For
\[
\exists y\in\R,\ \forall x\in\R,\ x+y=0,
\]
you must choose $y$ first. Then the adversary gets to choose $x$. Whatever $y$ you choose, the adversary can choose $x=1-y$, and then $x+y=1$, not $0$. You lose.

\begin{walkbox}
Ask yourself: in this theorem, who chooses first? Do I get to choose an object after seeing the input, or must I produce one object that works for all inputs?
\end{walkbox}

\newpage
\section{Negation: Turning Claims Inside Out}

Negation is not just putting ``not'' at the front. Good negation preserves the exact structure while reversing the truth.

\subsection{The basic laws}

\[
\neg(\forall x\in D,\ P(x)) \Longleftrightarrow \exists x\in D \text{ such that } \neg P(x).
\]

\[
\neg(\exists x\in D,\ P(x)) \Longleftrightarrow \forall x\in D,\ \neg P(x).
\]

If someone says, ``Every bridge is open,'' one closed bridge disproves it. If someone says, ``There is an open bridge,'' you must show every bridge is closed to disprove it.

\subsection{Negating implication}

The statement
\[
P\Rightarrow Q
\]
is false exactly when $P$ is true and $Q$ is false.

So:
\[
\neg(P\Rightarrow Q) \Longleftrightarrow P\land \neg Q.
\]

This matters when you test the boundary of a theorem.

Claim:

\begin{quote}
If a function is continuous, then it is differentiable.
\end{quote}

To disprove it, you need a function that is continuous and not differentiable. The absolute value function $f(x)=|x|$ at $0$ is the classic first example.

\subsection{Nested negations}

Negate:
\[
\forall \eps>0,\ \exists N\in\N \text{ such that } \forall n\geq N,\ |a_n-L|<\eps.
\]

Step by step:
\[
\exists \eps>0 \text{ such that } \forall N\in\N,\ \exists n\geq N \text{ such that } |a_n-L|\geq \eps.
\]

This is the negation of sequence convergence. In words: there is some fixed tolerance $\eps$ such that no matter how far out you go, the sequence still eventually escapes that tolerance.

That sentence is not just logic practice. It is analysis beginning to breathe.

\subsection{Exercises}

Negate the following.

\begin{enumerate}
\item Every integer is even.
\item There exists a real number $x$ such that $x^2<0$.
\item For every $x\in A$, there exists $y\in B$ such that $R(x,y)$.
\item If $n^2$ is even, then $n$ is even.
\item Every continuous function on $\R$ is bounded.
\end{enumerate}

\newpage
\section{If-Then: The Spine of Theorems}

Many theorems have the form:
\[
P\Rightarrow Q.
\]

The hypothesis $P$ is what you may assume. The conclusion $Q$ is what you must prove.

\subsection{Direct proof}

The direct proof is the plain road:

\begin{enumerate}
\item Assume $P$.
\item Use definitions and known facts.
\item Derive $Q$.
\end{enumerate}

\begin{theorem}
If $n$ is divisible by $6$, then $n$ is divisible by $3$.
\end{theorem}

\begin{proof}
Suppose $n$ is divisible by $6$. Then there exists $k\in\Z$ such that $n=6k$. Since $6k=3(2k)$ and $2k\in\Z$, there exists an integer $m=2k$ such that $n=3m$. Therefore $n$ is divisible by $3$.
\end{proof}

\subsection{What happened?}

The proof opened the definition of divisibility by $6$ and rearranged it into the definition of divisibility by $3$.

That is a large fraction of beginner proof writing.

\begin{focusbox}
When stuck, ask: what does the hypothesis mean by definition, and what does the conclusion require by definition?
\end{focusbox}

\subsection{Converse}

The converse of $P\Rightarrow Q$ is $Q\Rightarrow P$.

The theorem ``if $n$ is divisible by $6$, then $n$ is divisible by $3$'' is true. Its converse says ``if $n$ is divisible by $3$, then $n$ is divisible by $6$.'' False. Example: $n=3$.

Many false proofs accidentally prove or assume the converse.

\newpage
\section{Definitions Are Machines}

Definitions are where rigor begins. A definition tells you exactly what must be shown and exactly what may be used.

\begin{definition}
An integer $n$ is even if there exists $k\in\Z$ such that $n=2k$.
\end{definition}

\begin{definition}
An integer $n$ is odd if there exists $k\in\Z$ such that $n=2k+1$.
\end{definition}

\begin{definition}
For integers $a,b$, we say $a$ divides $b$, written $a\mid b$, if there exists $k\in\Z$ such that $b=ak$.
\end{definition}

\subsection{Opening and closing definitions}

If you know $n$ is even, you may open the definition:
\[
n=2k \text{ for some } k\in\Z.
\]

If you want to prove $n$ is even, you must close the definition: find an integer $k$ such that $n=2k$.

\subsection{Worked proof: even plus even}

\begin{theorem}
The sum of two even integers is even.
\end{theorem}

\begin{proof}
Let $a,b\in\Z$ be even. Since $a$ is even, there exists $m\in\Z$ such that $a=2m$. Since $b$ is even, there exists $n\in\Z$ such that $b=2n$. Then
\[
a+b=2m+2n=2(m+n).
\]
Since $m+n\in\Z$, the number $a+b$ has the form $2k$ for some integer $k$. Therefore $a+b$ is even.
\end{proof}

\subsection{Worked proof: divisibility is transitive}

\begin{theorem}
If $a,b,c\in\Z$, $a\mid b$, and $b\mid c$, then $a\mid c$.
\end{theorem}

\begin{proof}
Suppose $a\mid b$ and $b\mid c$. Since $a\mid b$, there exists $m\in\Z$ such that $b=am$. Since $b\mid c$, there exists $n\in\Z$ such that $c=bn$. Substituting $b=am$ into $c=bn$, we get
\[
c=(am)n=a(mn).
\]
Since $mn\in\Z$, this shows $a\mid c$.
\end{proof}

\begin{walkbox}
Pick a definition you know. Ask: if I assume it, what algebraic or logical object do I get? If I need to prove it, what object must I produce?
\end{walkbox}

\newpage
\section{The Three First Proof Moves}

At the start, you do not need a giant toolbox. You need a few reliable moves.

\subsection{Move 1: Direct proof}

Use it when the hypothesis naturally gives the conclusion.

\[
P\Rightarrow Q.
\]

Assume $P$. Prove $Q$.

\subsection{Move 2: Contrapositive}

The contrapositive of
\[
P\Rightarrow Q
\]
is
\[
\neg Q\Rightarrow \neg P.
\]

These are logically equivalent.

Use contrapositive when $\neg Q$ gives more concrete information than $P$.

\begin{theorem}
If $n^2$ is even, then $n$ is even.
\end{theorem}

\begin{proof}
We prove the contrapositive: if $n$ is not even, then $n^2$ is not even. For integers, not even means odd. Suppose $n$ is odd. Then $n=2k+1$ for some $k\in\Z$. Then
\[
n^2=(2k+1)^2=4k^2+4k+1=2(2k^2+2k)+1.
\]
Since $2k^2+2k\in\Z$, $n^2$ is odd. Therefore $n^2$ is not even. By contrapositive, if $n^2$ is even, then $n$ is even.
\end{proof}

\subsection{Move 3: Contradiction}

To prove $P$, assume $\neg P$ and derive an impossibility.

Contradiction is powerful, but beginners sometimes use it when direct proof would be cleaner. It is a sharp knife. Use it with care.

\begin{theorem}
$\sqrt2$ is irrational.
\end{theorem}

\begin{proof}
Suppose, for contradiction, that $\sqrt2$ is rational. Then $\sqrt2=a/b$ for integers $a,b$ with $b\neq0$, and we may choose the fraction in lowest terms. Squaring gives
\[
2=\frac{a^2}{b^2},
\]
so
\[
a^2=2b^2.
\]
Thus $a^2$ is even, so $a$ is even. Write $a=2k$. Then
\[
(2k)^2=2b^2,
\]
so
\[
4k^2=2b^2,
\]
and hence
\[
b^2=2k^2.
\]
Thus $b^2$ is even, so $b$ is even. Therefore both $a$ and $b$ are even, contradicting the assumption that $a/b$ was in lowest terms. Hence $\sqrt2$ is irrational.
\end{proof}

\newpage
\section{Examples and Counterexamples}

Examples are not decorations. They are how you test whether your understanding has touched the ground.

\subsection{Example}

A function $f:\R\to\R$ that is injective:
\[
f(x)=2x+1.
\]

If $f(a)=f(b)$, then $2a+1=2b+1$, so $a=b$.

\subsection{Non-example}

A function $f:\R\to\R$ that is not injective:
\[
f(x)=x^2.
\]

Because $f(1)=1=f(-1)$ but $1\neq -1$.

\subsection{Counterexample}

A counterexample is an example that disproves a universal statement.

Claim:
\[
\forall x,y\in\R,\ x^2=y^2\Rightarrow x=y.
\]

Counterexample: $x=1$, $y=-1$. Then $x^2=y^2$, but $x\neq y$.

\subsection{The useful cruelty of counterexamples}

Counterexamples are not mean. They are honest. A theorem with no boundary is usually not a theorem; it is a guess.

When a theorem is false, a counterexample shows where it breaks. When a theorem is true, failed counterexamples teach you which assumptions are secretly holding the roof up.

\subsection{Exercises}

Find counterexamples.

\begin{enumerate}
\item Every function $f:\R\to\R$ is surjective.
\item If $ab=0$, then $a=0$ and $b=0$.
\item If $A\cup B=A\cup C$, then $B=C$.
\item Every nonempty subset of $\R$ has a smallest element.
\item Every continuous function is differentiable.
\end{enumerate}

\newpage
\section{Sets: The Grammar of Modern Mathematics}

A set is a collection of objects. The objects are called elements.

\[
x\in A
\]
means $x$ is an element of $A$.

\[
x\notin A
\]
means $x$ is not an element of $A$.

Sets look innocent, but they became the common language of modern mathematics because they let us talk about collections of numbers, functions, points, symmetries, shapes, solutions, and open neighborhoods with one grammar.

\subsection{Basic operations}

\[
A\cup B=\{x:x\in A \text{ or } x\in B\}.
\]

\[
A\cap B=\{x:x\in A \text{ and } x\in B\}.
\]

\[
A\setminus B=\{x:x\in A \text{ and } x\notin B\}.
\]

\subsection{Subset}

\begin{definition}
$A\subseteq B$ means every element of $A$ is an element of $B$.
\end{definition}

To prove $A\subseteq B$:

\begin{quote}
Let $x\in A$ be arbitrary. Show $x\in B$.
\end{quote}

To prove $A=B$:

\begin{quote}
Prove $A\subseteq B$ and $B\subseteq A$.
\end{quote}

\subsection{Worked proof}

\begin{theorem}
For any sets $A$ and $B$,
\[
A\setminus B=A\cap B^c
\]
where $B^c$ is the complement of $B$ in the ambient universe.
\end{theorem}

\begin{proof}
Let $x$ be arbitrary. Then
\[
x\in A\setminus B
\]
if and only if $x\in A$ and $x\notin B$. But $x\notin B$ means $x\in B^c$. Therefore $x\in A$ and $x\in B^c$, which means $x\in A\cap B^c$. Thus $A\setminus B=A\cap B^c$.
\end{proof}

\subsection{Historical note}

Set theory became central in the nineteenth century, especially through Cantor's work on infinity. Before Cantor, infinity often behaved like smoke in the room: useful, present, but hard to hold. Set theory gave mathematicians a way to compare infinite collections with precision.

\newpage
\section{Set Proofs: Element Chasing}

Set proofs often reduce to following one element through definitions.

\begin{theorem}
For all sets $A,B,C$,
\[
A\cap(B\cup C)=(A\cap B)\cup(A\cap C).
\]
\end{theorem}

\begin{proof}
We prove both inclusions.

First suppose $x\in A\cap(B\cup C)$. Then $x\in A$ and $x\in B\cup C$. Since $x\in B\cup C$, either $x\in B$ or $x\in C$.

If $x\in B$, then $x\in A\cap B$, so $x\in (A\cap B)\cup(A\cap C)$.

If $x\in C$, then $x\in A\cap C$, so $x\in (A\cap B)\cup(A\cap C)$.

So $A\cap(B\cup C)\subseteq (A\cap B)\cup(A\cap C)$.

Now suppose $x\in (A\cap B)\cup(A\cap C)$. Then either $x\in A\cap B$ or $x\in A\cap C$.

If $x\in A\cap B$, then $x\in A$ and $x\in B$. Hence $x\in B\cup C$, so $x\in A\cap(B\cup C)$.

If $x\in A\cap C$, then $x\in A$ and $x\in C$. Hence $x\in B\cup C$, so $x\in A\cap(B\cup C)$.

So $(A\cap B)\cup(A\cap C)\subseteq A\cap(B\cup C)$. Therefore the sets are equal.
\end{proof}

\subsection{What to notice}

The proof did not draw a Venn diagram. Venn diagrams are useful for intuition, but they are not proof enough for arbitrary sets. The rigorous proof follows membership.

\begin{focusbox}
Set identities are where proof writing becomes grammar practice. Let $x$ live on one side. Translate definitions until it lives on the other.
\end{focusbox}

\subsection{Exercises}

Prove:

\begin{enumerate}
\item $A\cap B\subseteq A$.
\item $A\subseteq A\cup B$.
\item $A\setminus(B\cup C)=(A\setminus B)\cap(A\setminus C)$.
\item $A\setminus(B\cap C)=(A\setminus B)\cup(A\setminus C)$.
\end{enumerate}

\newpage
\section{Functions: More Than Formulas}

A function is often introduced as a formula, but a function is not just a formula.

\begin{definition}
A function $f:A\to B$ assigns to every element of $A$ exactly one element of $B$.
\end{definition}

The set $A$ is the domain. The set $B$ is the codomain.

The image is:
\[
\operatorname{im}(f)=\{f(a):a\in A\}.
\]

\subsection{Why codomain matters}

The formula $f(x)=x^2$ can define many different functions:

\[
f:\R\to\R,\qquad f:\R\to[0,\infty),\qquad f:\N\to\N.
\]

Same formula. Different function, because the domain and codomain are part of the data.

\subsection{Injective}

\begin{definition}
$f:A\to B$ is injective if for all $a_1,a_2\in A$,
\[
f(a_1)=f(a_2)\Rightarrow a_1=a_2.
\]
\end{definition}

Injective means no merging. If two outputs are equal, the inputs must have been equal.

\subsection{Surjective}

\begin{definition}
$f:A\to B$ is surjective if for every $b\in B$, there exists $a\in A$ such that $f(a)=b$.
\end{definition}

Surjective means every target in the codomain is hit.

\subsection{Bijective}

\begin{definition}
$f$ is bijective if it is both injective and surjective.
\end{definition}

Bijections are mathematical pairings. They say two sets have the same structure of elements, at least from the viewpoint of counting.

\newpage
\section{Function Proofs}

\begin{theorem}
If $f:A\to B$ and $g:B\to C$ are injective, then $g\circ f:A\to C$ is injective.
\end{theorem}

\begin{proof}
Let $a_1,a_2\in A$ and suppose $(g\circ f)(a_1)=(g\circ f)(a_2)$. Then
\[
g(f(a_1))=g(f(a_2)).
\]
Since $g$ is injective, $f(a_1)=f(a_2)$. Since $f$ is injective, $a_1=a_2$. Therefore $g\circ f$ is injective.
\end{proof}

\subsection{Inverse images}

If $f:A\to B$ and $S\subseteq B$, the inverse image of $S$ is:
\[
f^{-1}(S)=\{a\in A:f(a)\in S\}.
\]

This notation does not require $f$ to have an inverse function. It is a set operation.

\begin{theorem}
For $f:A\to B$ and $S,T\subseteq B$,
\[
f^{-1}(S\cap T)=f^{-1}(S)\cap f^{-1}(T).
\]
\end{theorem}

\begin{proof}
Let $a\in A$. Then
\[
a\in f^{-1}(S\cap T)
\]
if and only if
\[
f(a)\in S\cap T.
\]
This holds if and only if $f(a)\in S$ and $f(a)\in T$. That holds if and only if $a\in f^{-1}(S)$ and $a\in f^{-1}(T)$. Therefore $a\in f^{-1}(S)\cap f^{-1}(T)$. Hence the sets are equal.
\end{proof}

\subsection{Why inverse images matter later}

Topology defines continuity by inverse images:

\[
f \text{ is continuous if } f^{-1}(U) \text{ is open whenever } U \text{ is open.}
\]

That sentence is not beginner fluff. It is a bridge to topology. The habit of understanding $f^{-1}(S)$ now pays off later.

\newpage
\section{Relations: When Objects Count as Connected}

A relation on a set $A$ is a rule that tells whether two elements of $A$ are related.

Examples:

\begin{itemize}
\item $a=b$ on $\Z$.
\item $a\leq b$ on $\R$.
\item $a\mid b$ on $\Z$.
\item $a\equiv b\pmod n$ on $\Z$.
\end{itemize}

\subsection{Equivalence relations}

\begin{definition}
A relation $\sim$ on $A$ is an equivalence relation if it is reflexive, symmetric, and transitive.
\end{definition}

Reflexive:
\[
a\sim a.
\]

Symmetric:
\[
a\sim b\Rightarrow b\sim a.
\]

Transitive:
\[
a\sim b\text{ and }b\sim c\Rightarrow a\sim c.
\]

\subsection{Congruence modulo $n$}

Define $a\equiv b\pmod n$ if $n\mid(a-b)$.

\begin{theorem}
Congruence modulo $n$ is an equivalence relation on $\Z$.
\end{theorem}

\begin{proof}
Reflexive: For any $a\in\Z$, $a-a=0$, and $n\mid0$, so $a\equiv a\pmod n$.

Symmetric: Suppose $a\equiv b\pmod n$. Then $n\mid(a-b)$, so $a-b=nk$ for some $k\in\Z$. Then $b-a=-nk=n(-k)$, so $n\mid(b-a)$. Hence $b\equiv a\pmod n$.

Transitive: Suppose $a\equiv b\pmod n$ and $b\equiv c\pmod n$. Then $n\mid(a-b)$ and $n\mid(b-c)$. So $a-b=nm$ and $b-c=n\ell$ for integers $m,\ell$. Adding,
\[
a-c=(a-b)+(b-c)=nm+n\ell=n(m+\ell).
\]
Thus $n\mid(a-c)$, so $a\equiv c\pmod n$.
\end{proof}

\subsection{Why relations matter later}

Equivalence relations create quotient structures. Quotients are everywhere: modular arithmetic, quotient groups, quotient rings, projective spaces, manifolds built by gluing, and more. A quotient says: we will treat these different-looking things as the same under a chosen rule.

\newpage
\section{Induction: Infinite Proof With Finite Labor}

Induction is one of the first places mathematics feels like a magic trick until you see the mechanism.

\begin{theorem}[Principle of induction]
Suppose $P(n)$ is a statement for each integer $n\geq 1$. If $P(1)$ is true, and for every $k\geq1$, $P(k)\Rightarrow P(k+1)$, then $P(n)$ is true for all $n\geq1$.
\end{theorem}

The image is dominoes. Knock down the first. Show each knocks down the next. Then all fall.

\subsection{Worked proof}

\begin{theorem}
For every $n\geq1$,
\[
1+2+\cdots+n=\frac{n(n+1)}{2}.
\]
\end{theorem}

\begin{proof}
Base case: For $n=1$, the left side is $1$ and the right side is $\frac{1(2)}2=1$. So the statement holds for $n=1$.

Induction step: Assume for some $k\geq1$ that
\[
1+2+\cdots+k=\frac{k(k+1)}2.
\]
We must prove
\[
1+2+\cdots+k+(k+1)=\frac{(k+1)(k+2)}2.
\]
Using the induction hypothesis,
\[
1+2+\cdots+k+(k+1)=\frac{k(k+1)}2+(k+1).
\]
Then
\[
\frac{k(k+1)}2+(k+1)=\frac{k(k+1)+2(k+1)}2
=\frac{(k+1)(k+2)}2.
\]
So the statement holds for $k+1$. By induction, it holds for all $n\geq1$.
\end{proof}

\subsection{Common induction bug}

Do not assume what you are trying to prove for $k+1$. The induction hypothesis gives you $P(k)$, not $P(k+1)$.

\begin{focusbox}
In induction, always label the base case, induction hypothesis, and induction step.
\end{focusbox}

\newpage
\section{Strong Induction and Factorization}

Strong induction lets you assume all earlier cases, not just the immediately previous one.

\begin{theorem}[Strong induction]
Suppose $P(n)$ is a statement for integers $n\geq n_0$. If $P(n_0)$ is true, and for every $k\geq n_0$, the truth of $P(n_0),P(n_0+1),\ldots,P(k)$ implies $P(k+1)$, then $P(n)$ is true for all $n\geq n_0$.
\end{theorem}

\subsection{Prime factorization existence}

\begin{theorem}
Every integer $n\geq2$ is a product of primes.
\end{theorem}

\begin{proof}
We use strong induction on $n$.

Base case: $2$ is prime, so it is a product of primes.

Induction step: Suppose every integer $m$ with $2\leq m\leq k$ is a product of primes. We prove $k+1$ is a product of primes.

If $k+1$ is prime, then it is already a product of primes.

If $k+1$ is not prime, then $k+1=ab$ for integers $a,b$ with $2\leq a\leq k$ and $2\leq b\leq k$. By the induction hypothesis, both $a$ and $b$ are products of primes. Therefore $ab=k+1$ is a product of primes.

Thus by strong induction, every integer $n\geq2$ is a product of primes.
\end{proof}

\subsection{Why this matters}

This proof is a doorway to number theory. Later, unique factorization becomes a central theme. In ordinary integers, factorization behaves well. In other number systems, it can fail. That failure forced mathematicians to invent ideals, which later became essential in algebraic geometry.

\begin{walkbox}
Think about software dependencies. A package can depend on smaller packages, which depend on smaller packages, until you reach primitives. Strong induction is the proof version of trusting all smaller dependencies.
\end{walkbox}

\newpage
\section{Equality, Identity, and Isomorphism}

Beginners often think mathematics only cares about equality. Later mathematics cares just as much about sameness of structure.

\subsection{Equality}

Equality means two expressions name the same object.

\[
2+3=5.
\]

\subsection{Isomorphism}

An isomorphism means two structures may have different labels but the same pattern.

For example, the even integers under addition behave like the integers under addition:
\[
\Z \to 2\Z,\qquad n\mapsto 2n.
\]

This map preserves addition:
\[
2(a+b)=2a+2b.
\]

The labels changed. The additive structure did not.

\subsection{Why this matters}

Abstract algebra is full of this idea. Topology has homeomorphism. Differential geometry has diffeomorphism. Category theory makes structure-preserving maps the main characters.

The deep habit: do not ask only whether two things are equal. Ask what structure is preserved by the map between them.

\newpage
\section{Infinity: Countable and Uncountable}

Infinity is where mathematics learned humility.

The natural numbers are infinite:
\[
\N=\{1,2,3,\ldots\}.
\]

The even positive integers are also infinite:
\[
\{2,4,6,\ldots\}.
\]

There are ``half as many'' evens in an everyday sense, but the map
\[
n\mapsto 2n
\]
pairs every natural number with exactly one even positive integer. So the two sets have the same cardinality.

\begin{definition}
A set is countably infinite if its elements can be listed in a sequence indexed by $\N$.
\end{definition}

\subsection{The integers are countable}

List them:
\[
0,1,-1,2,-2,3,-3,\ldots
\]

Every integer appears eventually.

\subsection{The rationals are countable}

Fractions seem denser. Between any two rationals, there is another rational. Still, $\Q$ is countable.

Imagine a grid of fractions $p/q$ with $p\in\Z$ and $q\in\N$. Sweep through the grid by increasing $|p|+q$, skipping duplicates like $2/4=1/2$. Every rational appears eventually.

\subsection{The reals are uncountable}

Cantor's diagonal argument shows no list of real numbers in $(0,1)$ can be complete. Given any proposed list of decimals, build a new decimal that differs from the first number in the first digit, the second number in the second digit, and so on. The new number differs from every listed number, so it was missing.

Infinity has sizes.

\subsection{Why this matters later}

Analysis lives among uncountable sets. Topology studies spaces where infinity can behave subtly. Algebraic geometry often studies solution sets with infinitely many points controlled by finite equations.

\newpage
\section{First Epsilon Proofs}

Engineering calculus often treats limits as motion. A variable approaches. A graph settles. A value tends toward another value.

Real analysis asks for the contract.

\begin{definition}
A sequence $(a_n)$ converges to $L$ if for every $\eps>0$, there exists $N\in\N$ such that for all $n\geq N$,
\[
|a_n-L|<\eps.
\]
\end{definition}

Translation: no matter how small a tolerance band you choose around $L$, the sequence eventually stays inside it forever.

\subsection{Worked proof}

\begin{theorem}
\[
\lim_{n\to\infty}\frac1n=0.
\]
\end{theorem}

\begin{proof}
Let $\eps>0$ be arbitrary. We need to find $N\in\N$ such that for all $n\geq N$,
\[
\left|\frac1n-0\right|<\eps.
\]
Choose $N\in\N$ such that $N>1/\eps$. Then if $n\geq N$,
\[
0<\frac1n\leq\frac1N<\eps.
\]
Thus $\left|\frac1n-0\right|<\eps$. Therefore $1/n\to0$.
\end{proof}

\subsection{The secret}

You do not choose $\eps$. The challenger chooses $\eps$. You choose $N$ in response.

\begin{walkbox}
Think of $\eps$ as a tolerance in a sensor. Someone sets the tolerance. Your job is to say how far into the sequence you must go before the readings stay within tolerance forever.
\end{walkbox}

\newpage
\section{Continuity: Local Control}

The epsilon-delta definition of continuity looks intimidating because it has nested quantifiers. But the idea is physical:

\begin{quote}
Small changes in input force small changes in output.
\end{quote}

\begin{definition}
A function $f:\R\to\R$ is continuous at $c\in\R$ if for every $\eps>0$, there exists $\delta>0$ such that for all $x\in\R$, if $|x-c|<\delta$, then $|f(x)-f(c)|<\eps$.
\end{definition}

\subsection{Worked proof}

\begin{theorem}
The function $f(x)=3x+2$ is continuous at every $c\in\R$.
\end{theorem}

\begin{proof}
Let $c\in\R$ and let $\eps>0$ be arbitrary. We need $\delta>0$ such that $|x-c|<\delta$ implies $|f(x)-f(c)|<\eps$.

Now
\[
|f(x)-f(c)|=|(3x+2)-(3c+2)|=|3x-3c|=3|x-c|.
\]
Choose $\delta=\eps/3$. If $|x-c|<\delta$, then
\[
|f(x)-f(c)|=3|x-c|<3\delta=3(\eps/3)=\eps.
\]
Therefore $f$ is continuous at $c$.
\end{proof}

\subsection{Connection to topology}

Topology later replaces epsilon-delta with open sets. Instead of measuring distances directly, it says a function is continuous when inverse images of open sets are open. That is a huge abstraction, but it is not alien. It is the same idea: local neighborhoods are respected.

\newpage
\section{Algebra Begins With Symmetry}

Abstract algebra can look like someone removed all the numbers and left only laws. That is partly right. Algebra asks what can be known when the objects change but the operation rules remain.

\begin{definition}
A group is a set $G$ with an operation $\cdot$ such that:
\begin{enumerate}
\item closure: if $a,b\in G$, then $a\cdot b\in G$,
\item associativity: $(a\cdot b)\cdot c=a\cdot(b\cdot c)$,
\item identity: there exists $e\in G$ such that $e\cdot a=a\cdot e=a$ for all $a\in G$,
\item inverses: for each $a\in G$, there exists $a^{-1}\in G$ such that $a\cdot a^{-1}=a^{-1}\cdot a=e$.
\end{enumerate}
\end{definition}

Examples:

\begin{itemize}
\item $\Z$ under addition.
\item Nonzero real numbers under multiplication.
\item Symmetries of a square under composition.
\item Invertible matrices under multiplication.
\end{itemize}

\subsection{Worked proof: identity is unique}

\begin{theorem}
The identity element of a group is unique.
\end{theorem}

\begin{proof}
Suppose $e$ and $e'$ are both identity elements in a group $G$. Since $e$ is an identity, $e\cdot e'=e'$. Since $e'$ is an identity, $e\cdot e'=e$. Therefore $e=e'$. Hence the identity element is unique.
\end{proof}

\subsection{What to notice}

The proof did not compute. It used only the definition. That is abstract algebra in miniature.

\newpage
\section{Topology Begins With Nearness Without Rulers}

Topology studies continuity and shape in a way that does not always require measuring distance.

\begin{definition}
A topology on a set $X$ is a collection $\mathcal{T}$ of subsets of $X$, called open sets, such that:
\begin{enumerate}
\item $\emptyset$ and $X$ are in $\mathcal{T}$,
\item arbitrary unions of sets in $\mathcal{T}$ are in $\mathcal{T}$,
\item finite intersections of sets in $\mathcal{T}$ are in $\mathcal{T}$.
\end{enumerate}
\end{definition}

This definition looks sparse. But it captures enough structure to talk about continuity.

\begin{definition}
If $X$ and $Y$ are topological spaces, a function $f:X\to Y$ is continuous if for every open set $U\subseteq Y$, the inverse image $f^{-1}(U)$ is open in $X$.
\end{definition}

Now the earlier work on inverse images returns.

\subsection{Physical intuition}

Topology cares about neighborhoods. If a point lies in an open set, it has some room around it. The exact distance may not matter. What matters is the pattern of which neighborhoods exist.

This is why a coffee mug and a donut are famously similar in topology: each has one hole. Topology ignores bending and stretching but notices tearing and gluing.

\subsection{Warning}

The coffee mug story is memorable, but do not let it replace definitions. The definitions are the engine. The story is the spark plug.

\newpage
\section{Differential Geometry: Calculus on Curved Spaces}

Differential geometry begins when calculus leaves flat Euclidean space.

A curve can be studied with derivatives. A surface can be studied with tangent planes. A manifold is the general idea: a space that may be globally complicated but locally resembles $\R^n$.

\begin{definition}
Informally, an $n$-dimensional manifold is a space that near each point looks like $\R^n$.
\end{definition}

The phrase ``looks like'' becomes precise using charts.

\begin{definition}
A chart is a coordinate map from part of a manifold to an open subset of $\R^n$.
\end{definition}

\subsection{Connection to previous material}

To understand manifolds, you need:

\begin{itemize}
\item functions and inverse images,
\item open sets and continuity,
\item linear algebra,
\item multivariable calculus,
\item proof comfort.
\end{itemize}

The tangent space at a point is a linear approximation to a curved object. The derivative becomes a map between tangent spaces.

\subsection{Physical connection}

General relativity uses differential geometry because spacetime is modeled as a manifold with curvature. The mathematics of curved spaces becomes physics when curvature tells matter how to move and matter tells curvature how to change.

\newpage
\section{Algebraic Geometry: Equations Become Shapes}

Algebraic geometry begins with a simple observation:

\begin{quote}
Polynomial equations define geometric sets.
\end{quote}

For example:
\[
x^2+y^2=1
\]
defines a circle in the plane.

\[
y^2=x^3-x
\]
defines an elliptic curve.

\subsection{Affine varieties}

\begin{definition}
Given polynomials $f_1,\ldots,f_m$ in variables $x_1,\ldots,x_n$ over a field $k$, their common zero set is
\[
V(f_1,\ldots,f_m)=\{a\in k^n:f_i(a)=0\text{ for all }i\}.
\]
\end{definition}

This is a first version of an affine variety.

\subsection{Ideals}

Instead of studying one polynomial at a time, algebra studies collections of polynomials closed under addition and multiplication by arbitrary polynomials. These are ideals.

\begin{definition}
An ideal $I$ in a ring $R$ is a subset such that:
\begin{enumerate}
\item if $a,b\in I$, then $a+b\in I$,
\item if $r\in R$ and $a\in I$, then $ra\in I$.
\end{enumerate}
\end{definition}

Ideals are algebraic machines for encoding equations.

\subsection{The bridge}

Geometry asks about shapes. Algebra asks about rings and ideals. Algebraic geometry learns to translate between them.

That is why the path here matters. Sets, functions, relations, quotients, topology, rings, ideals: these are not random hurdles. They are the grammar of the later language.

\newpage
\section{How to Read a Definition}

When you meet a definition, do not only underline it. Work it.

\subsection{The five-part definition routine}

For every important definition, write:

\begin{enumerate}
\item the formal definition,
\item a plain-English version,
\item two examples,
\item one near miss,
\item one theorem where the definition is used.
\end{enumerate}

\subsection{Example: injective}

Formal:
\[
f:A\to B \text{ is injective if } f(a_1)=f(a_2)\Rightarrow a_1=a_2.
\]

Plain English: the function does not collapse two different inputs to one output.

Examples:

\[
f:\R\to\R,\quad f(x)=2x+1.
\]

\[
g:\N\to\N,\quad g(n)=n+10.
\]

Near miss:

\[
h:\R\to\R,\quad h(x)=x^2.
\]

Theorem:

If $f$ and $g$ are injective, then $g\circ f$ is injective.

\begin{focusbox}
Definitions become memory when you can produce examples and non-examples without looking.
\end{focusbox}

\newpage
\section{How to Read a Theorem}

A theorem should be interrogated before it is proved.

Take:

\begin{quote}
If $f:A\to B$ and $g:B\to C$ are injective, then $g\circ f:A\to C$ is injective.
\end{quote}

Ask:

\begin{enumerate}
\item What are the objects? Sets $A,B,C$ and functions $f,g$.
\item What are the assumptions? $f$ and $g$ are injective.
\item What must be shown? $g\circ f$ is injective.
\item What does the conclusion mean by definition? If $(g\circ f)(a_1)=(g\circ f)(a_2)$, then $a_1=a_2$.
\item Which assumptions can help? Injectivity of $g$ first, then injectivity of $f$.
\end{enumerate}

\subsection{The theorem-reading template}

\begin{verbatim}
Objects:
Hypotheses:
Conclusion:
Definitions involved:
Likely proof strategy:
Possible counterexample if an assumption is removed:
\end{verbatim}

This template slows you down in the right place.

\newpage
\section{How to Write a Proof}

A proof is a piece of communication. It should show why the claim follows.

\subsection{Proof skeleton}

Before writing full sentences, write the skeleton.

For a universal implication:

\begin{quote}
Let [object] be arbitrary. Suppose [hypothesis]. We must show [conclusion]. By definition, [unpack]. Therefore [derive]. Hence [conclusion].
\end{quote}

For set equality:

\begin{quote}
We prove both inclusions. First let $x\in$ left side. ... Thus left $\subseteq$ right. Now let $x\in$ right side. ... Thus right $\subseteq$ left.
\end{quote}

For existence:

\begin{quote}
Choose [witness]. Verify [property]. Therefore such an object exists.
\end{quote}

\subsection{Sentences to avoid}

Avoid:

\begin{itemize}
\item ``Clearly''
\item ``Obviously''
\item ``It is easy to see''
\item ``By math''
\end{itemize}

These are sometimes harmless in expert writing, but early on they often hide the one step that needs attention.

\subsection{A better replacement}

Instead of:

\begin{quote}
Clearly, $a+b$ is even.
\end{quote}

Write:

\begin{quote}
Since $a+b=2(m+n)$ and $m+n\in\Z$, the number $a+b$ is even by definition.
\end{quote}

\newpage
\section{Proof Debugging}

Proof debugging is the habit that turns confusion into progress.

\subsection{The first broken line}

When a proof is wrong, do not rewrite the whole thing immediately. Find the first line that does not follow.

\subsection{Gap types}

\begin{longtable}{p{0.22\linewidth}p{0.32\linewidth}p{0.35\linewidth}}
\toprule
Gap & What it looks like & Repair \\
\midrule
Definition gap & A term is used but not unpacked & Write the definition \\
Quantifier gap & A variable appears without scope & Introduce it with domain and role \\
Existence gap & An object is claimed without witness & Construct or cite existence \\
Case gap & One possibility is handled & List all cases \\
Implication gap & A conclusion jumps too far & Insert a lemma or algebra step \\
Converse error & $P\Rightarrow Q$ treated as $Q\Rightarrow P$ & Check direction \\
\bottomrule
\end{longtable}

\subsection{Bad proof, repaired}

Claim: If $n^2$ is even, then $n$ is even.

Bad proof:

\begin{quote}
Since $n^2$ is even, $n^2=2k$. Taking square roots, $n=\sqrt{2k}$, so $n$ is even.
\end{quote}

Problem: being $\sqrt{2k}$ does not mean being twice an integer. The proof has an implication gap.

Repair: use contrapositive, as shown earlier.

\newpage
\section{Exercises for Walking}

These are not calculation exercises. They are thinking exercises for when you are away from paper.

\begin{enumerate}
\item What does it mean to prove a statement for an arbitrary object?
\item Why does one counterexample disprove a universal claim?
\item What is the difference between a function's codomain and image?
\item Why does injective mean ``no merging''?
\item Why does surjective mean ``full coverage''?
\item What is a quotient trying to forget?
\item Why is induction not circular?
\item In an epsilon proof, who chooses $\eps$?
\item Why does topology care about inverse images?
\item What does algebra preserve when it stops caring about particular numbers?
\end{enumerate}

Try answering in full sentences. If you cannot explain an idea while walking, you probably know its symbol better than its meaning.

\newpage
\section{A Four-Week Beginner-to-Velleman Launch}

\begin{longtable}{p{0.08\linewidth}p{0.24\linewidth}p{0.32\linewidth}p{0.26\linewidth}}
\toprule
Week & Theme & Work & Output \\
\midrule
1 & Sentences and quantifiers & Statements, predicates, domains, $\forall$, $\exists$, negation & 30 translations and 20 negations \\
2 & Definitions and direct proof & Even, odd, divides, subset, direct proof & 10 short proofs from definitions \\
3 & Sets and functions & Subsets, equality, injective, surjective, inverse image & 6 set proofs and 4 function proofs \\
4 & Relations and induction & Equivalence relations, modular arithmetic, induction & 3 relation proofs and 3 induction proofs \\
\bottomrule
\end{longtable}

\subsection{Daily rhythm}

\begin{enumerate}
\item Read 3 to 5 pages.
\item Copy one definition by hand.
\item Make two examples and one non-example.
\item Prove one small claim.
\item Debug yesterday's proof.
\item Review five spaced repetition prompts.
\end{enumerate}

\subsection{When to start Velleman}

Start Velleman when you can:

\begin{itemize}
\item negate nested quantifiers without panic,
\item prove simple parity and divisibility facts,
\item prove set equality by two inclusions,
\item prove injectivity statements from the definition,
\item write an induction proof with all parts labeled.
\end{itemize}

\newpage
\section{Milestone Proof Portfolio}

Before moving deeply into Velleman, write these ten proofs cleanly.

\begin{enumerate}
\item The sum of two even integers is even.
\item The product of two odd integers is odd.
\item If $a\mid b$ and $b\mid c$, then $a\mid c$.
\item If $n^2$ is even, then $n$ is even.
\item $\sqrt2$ is irrational.
\item $A\cap(B\cup C)=(A\cap B)\cup(A\cap C)$.
\item $f^{-1}(S\cap T)=f^{-1}(S)\cap f^{-1}(T)$.
\item The composition of injective functions is injective.
\item Congruence modulo $n$ is an equivalence relation.
\item $1+2+\cdots+n=n(n+1)/2$ for all $n\geq1$.
\end{enumerate}

For each proof, include:

\begin{itemize}
\item theorem statement,
\item definitions used,
\item proof strategy,
\item final proof,
\item one sentence explaining the key idea.
\end{itemize}

\begin{focusbox}
This portfolio is the bridge. If you can write these proofs cleanly, Velleman's first chapters will feel like organized terrain rather than fog.
\end{focusbox}

\newpage
\section{Closing: The Shape of the Road}

The road from engineering mathematics to pure mathematics is not a rejection of computation. It is a widening of what counts as understanding.

At first, proof feels slow because every little word matters. Then the slowness becomes sight. You begin to notice why a theorem needs each hypothesis. You begin to hear the difference between ``there exists'' and ``for every.'' You begin to see definitions as tools, examples as tests, counterexamples as boundary stones, and proofs as honest explanations.

From here, Velleman teaches the craft more systematically. Real analysis will sharpen your relationship with infinity and approximation. Abstract algebra will train you to recognize structure. Topology will separate continuity from measurement. Differential geometry will bring calculus onto curved spaces. Algebraic geometry will let equations and shapes speak to each other.

The beginning is humble: parse the sentence. Find the objects. Open the definitions. Move one justified step at a time.

That is not a small beginning. It is the beginning of doing mathematics.

\newpage
\appendix
\section{Appendix A: Notation You Should Be Able to Read}

\begin{longtable}{p{0.18\linewidth}p{0.34\linewidth}p{0.38\linewidth}}
\toprule
Notation & Meaning & What to ask yourself \\
\midrule
$x\in A$ & $x$ is an element of $A$ & What kind of objects live in $A$? \\
$A\subseteq B$ & Every element of $A$ is in $B$ & How would I prove this with an arbitrary $x$? \\
$A\cup B$ & Union & Is membership ``or''? \\
$A\cap B$ & Intersection & Is membership ``and''? \\
$A\setminus B$ & Difference & Is $x$ in $A$ but not in $B$? \\
$\pow(A)$ & Power set & What are the subsets of $A$? \\
$f:A\to B$ & Function from $A$ to $B$ & What are domain and codomain? \\
$f(S)$ & Image of $S$ & Which outputs come from inputs in $S$? \\
$f^{-1}(T)$ & Inverse image of $T$ & Which inputs land in $T$? \\
$\forall$ & For every & Who chooses the object? \\
$\exists$ & There exists & What witness proves it? \\
$P\Rightarrow Q$ & If $P$, then $Q$ & What is assumed and what is owed? \\
$P\Leftrightarrow Q$ & If and only if & Can I prove both directions? \\
$a\mid b$ & $a$ divides $b$ & Can I write $b=ak$? \\
$a\equiv b\pmod n$ & Same remainder mod $n$ & Is $n\mid(a-b)$? \\
$\epsilon,\delta$ & small positive tolerances & Which one is chosen first? \\
\bottomrule
\end{longtable}

\subsection{Notation drills}

Translate each into English.

\begin{enumerate}
\item $\forall x\in A,\ x\in B$.
\item $\exists x\in A$ such that $x\notin B$.
\item $f^{-1}(S\cup T)=f^{-1}(S)\cup f^{-1}(T)$.
\item $\forall \eps>0,\exists N\in\N,\forall n\geq N,\ |a_n-L|<\eps$.
\item $\forall a,b\in G,\ \phi(ab)=\phi(a)\phi(b)$.
\end{enumerate}

\newpage
\section{Appendix B: A Short History of Why Rigor Arrived}

Mathematics did not begin as a perfectly formal machine. It began with counting, measuring, astronomy, land, trade, architecture, music, and motion. Rigor arrived in waves whenever intuition became too powerful to trust without guardrails.

\subsection{Euclid and the axiomatic dream}

Euclid's \emph{Elements} organized geometry around definitions, postulates, and propositions. It was not modern formalism, but it established a model: start from agreed assumptions and reason carefully. A proof became a public object, something another mind could inspect.

\subsection{Calculus and the problem of motion}

Newton and Leibniz gave mathematics a language for change. Calculus worked spectacularly. It explained falling bodies, planetary motion, tangents, areas, and rates. But early calculus leaned on infinitesimals and limiting arguments that were powerful but not always clean.

The results were too useful to abandon, so mathematicians sharpened the foundations.

\subsection{Cauchy and Weierstrass}

The nineteenth century rebuilt calculus into analysis. Limits became epsilon-delta statements. Continuity became a quantified contract. Infinite series were no longer trusted just because they looked well-behaved.

That is why your work with quantifiers matters. The formal definition of a limit is not a hazing ritual. It is the historical repair of a brilliant machine.

\subsection{Cantor and infinity}

Cantor showed that infinity has sizes. The natural numbers, integers, and rationals are countable. The real numbers are not. This was shocking because it made infinity into something that could be compared, classified, and proved about.

\subsection{Noether and structure}

Emmy Noether helped push algebra away from individual equations and toward structures: groups, rings, fields, ideals. This shift is one reason modern mathematics feels abstract. Used well, abstraction compresses many concrete calculations into a single theorem that covers them all at once.

\subsection{Hilbert and formal clarity}

Hilbert championed axiomatic thinking. The idea was not to remove meaning from mathematics, but to clarify exactly what assumptions were being used. This habit still shapes modern proof writing.

\subsection{Grothendieck and the expansion of geometry}

Grothendieck transformed algebraic geometry by changing the level of abstraction. Instead of studying only solution sets of polynomials over familiar fields, he built a language where geometry could be studied through rings, schemes, and maps. That world is far away from this reader, but the early tools are already here: sets, functions, equivalence, quotients, topology, and algebra.

\begin{focusbox}
Rigor usually appears after intuition succeeds. First people find a powerful idea. Then they discover the idea needs a stronger language.
\end{focusbox}

\newpage
\section{Appendix C: Extra Proof Drills}

\subsection{Drill block 1: Definitions only}

Prove each using only definitions and basic algebra.

\begin{enumerate}
\item If $n$ is even, then $n^2$ is even.
\item If $m$ is even and $n$ is odd, then $mn$ is even.
\item If $a\mid b$ and $a\mid c$, then $a\mid(b+c)$.
\item If $a\mid b$, then $a\mid bc$ for every $c\in\Z$.
\item If $a$ and $b$ are rational, then $a+b$ is rational.
\item If $a$ and $b$ are rational, then $ab$ is rational.
\end{enumerate}

\subsection{Drill block 2: Counterexamples}

Disprove each false statement.

\begin{enumerate}
\item If $a\mid bc$, then $a\mid b$ or $a\mid c$.
\item If $x<y$, then $x^2<y^2$.
\item Every injective function is surjective.
\item Every surjective function is injective.
\item If $A\subseteq B\cup C$, then $A\subseteq B$ or $A\subseteq C$.
\item If $A\cap B=A\cap C$, then $B=C$.
\end{enumerate}

\subsection{Drill block 3: Proof strategy choice}

For each theorem, choose a likely strategy before proving it.

\begin{enumerate}
\item If $n^2$ is odd, then $n$ is odd.
\item There is no smallest positive real number.
\item The intersection of two subgroups of a group is a subgroup.
\item If $A\subseteq B$, then $\pow(A)\subseteq\pow(B)$.
\item If $f:A\to B$ is injective and $S,T\subseteq A$, then $f(S\cap T)=f(S)\cap f(T)$.
\end{enumerate}

\newpage
\section{Appendix D: Hints Without Full Solutions}

\subsection{Definitions only}

\begin{enumerate}
\item For $n^2$ even: write $n=2k$, then square.
\item For even times odd: write $m=2a$ and $n=2b+1$.
\item For divisibility over sums: write $b=ak$ and $c=a\ell$.
\item For $a\mid bc$: start from $b=ak$.
\item For rational sums: write $a=p/q$ and $b=r/s$ with nonzero denominators.
\item For rational products: multiply numerators and denominators.
\end{enumerate}

\subsection{Counterexamples}

\begin{enumerate}
\item Try $a=6$, $b=2$, $c=3$.
\item Try negative numbers.
\item Think of $f:\N\to\N$ given by $f(n)=n+1$.
\item Think of a map from a two-element set to a one-element set.
\item Let $A$ contain two points, one in $B$ and one in $C$.
\item Make $A$ empty.
\end{enumerate}

\subsection{Strategy choice}

\begin{enumerate}
\item Contrapositive is clean.
\item Contradiction is natural: assume a smallest positive real exists and halve it.
\item Direct proof from subgroup definition.
\item Let $S\in\pow(A)$ and prove $S\in\pow(B)$.
\item Prove both inclusions; injectivity is needed for the harder direction.
\end{enumerate}

\newpage
\section{Appendix E: Reading Velleman Without Getting Lost}

When you open Velleman, do not read it like a novel and do not read it like a formula sheet. Read it like a manual for a new instrument.

\subsection{Before each section}

\begin{enumerate}
\item Skim the section headings.
\item Write down the symbols you expect to see.
\item Predict what kind of proof method may appear.
\end{enumerate}

\subsection{During each section}

\begin{enumerate}
\item Copy every definition that seems central.
\item After each example proof, write the strategy in one line.
\item Mark any sentence where a variable is introduced.
\item Pause at every ``therefore'' and ask what justifies it.
\end{enumerate}

\subsection{After each section}

\begin{enumerate}
\item Rewrite one proof from memory.
\item Make one example and one non-example for each definition.
\item Do at least one exercise badly before looking for help.
\item Debug the bad attempt before writing a polished version.
\end{enumerate}

\subsection{What ``understanding'' should feel like}

Understanding does not mean every proof feels obvious. It means you can locate the difficulty. You can say: ``I understand the definitions, but I do not see why this witness works,'' or ``I see the strategy, but the induction step has an algebra gap.''

That is progress. Fog becomes a map when you can name where you are.

\newpage
\section{Appendix F: The One-Page Proof Debugger}

Use this every time a proof feels shaky.

\begin{enumerate}
\item State the claim exactly.
\item Circle all quantifiers.
\item List the objects and their domains.
\item Separate hypotheses from conclusion.
\item Expand every definition in the conclusion.
\item Expand useful definitions in the hypotheses.
\item Choose a strategy.
\item Write a skeleton.
\item Fill in one line at a time.
\item For each line, ask: why does this follow?
\item Check that arbitrary objects were really arbitrary.
\item Check that witnesses were actually constructed.
\item Check all cases.
\item Remove ``clearly'' unless the next sentence explains why.
\item Read the final proof aloud.
\end{enumerate}

\begin{focusbox}
The proof debugger is not remedial. Professionals debug proofs too. The difference is that they have better names for the failure modes.
\end{focusbox}

\newpage
\section{Appendix G: Final Readiness Exam}

Try this without notes. It is not a speed test.

\begin{enumerate}
\item Define statement, predicate, theorem, proof, counterexample.
\item Negate: $\forall x\in A,\exists y\in B,\ R(x,y)$.
\item Prove that if $n$ is odd, then $n^2$ is odd.
\item Prove that if $a\mid b$ and $a\mid c$, then $a\mid(b-c)$.
\item Prove $A\setminus(B\cup C)=(A\setminus B)\cap(A\setminus C)$.
\item Give a counterexample to: if $f(S\cap T)=f(S)\cap f(T)$, then $f$ is injective.
\item Prove that if $f$ and $g$ are surjective, then $g\circ f$ is surjective.
\item Prove congruence modulo $n$ is transitive.
\item Prove by induction that $2^n\geq n+1$ for all $n\geq1$.
\item Explain, in words, why $\Q$ is countable but $\R$ is not.
\end{enumerate}

\subsection{How to score yourself}

\begin{itemize}
\item If you cannot start: review definitions.
\item If you start but drift: review quantifiers and proof skeletons.
\item If your proof has gaps: use the debugger.
\item If your proof is correct but ugly: good. Polish comes after correctness.
\end{itemize}

\end{document}
